% Section 5 — The model: problem encoding, kernel embedding, game and objectives
\section{From data to algorithm: the InnovationZero model}
\label{sec:model}

The InnovationZero model is an \emph{algorithmic foundation model}: a network whose output is an
algorithm, and whose training signal is what that algorithm does on data it has not seen. This
section describes it at the level a reader needs to assess the claims, and no further at one
point: the language in which kernels are written is proprietary, so we describe it by the
guarantees it enforces rather than by its syntax (\cref{app:disclosure} draws the boundary
precisely). Four maps define the model:
\begin{equation}
z = E_\phi(D_{\mathrm{tr}}), \qquad h = S_\theta(z), \qquad a = G(h), \qquad c_\ell = \Lambda_\ell(a),
\label{eq:contract}
\end{equation}
an encoder $E_\phi$ from training rows to a \emph{problem embedding} $z$; the \emph{kernel
embedding} $S_\theta$, which turns $z$ into a representation $h$ of a solution and decodes it into
a typed kernel $a=G(h)$; and emitters $\Lambda_\ell$ that lower the kernel to source code in a
language $\ell$. After commitment, the emitted program is fitted and scored exactly as in
\cref{eq:commit}. \Cref{fig:pipeline} shows the flow.

\begin{figure}[t]
\centering
% The data-to-algorithm flow. The dark stage is the protected kernel embedding.
\begin{tikzpicture}[node distance=5.5mm and 4.2mm]
\node[izbox, minimum height=13mm, text width=21mm] (data) {Training rows\\[1pt]{\scriptsize target in last column;\\held-out rows never enter}};
\node[izstage, right=of data, minimum height=13mm, text width=20mm] (enc) {Problem encoder\\[1pt]{\scriptsize Hyperion v7\\transformer}};
\node[izbox, right=of enc, minimum height=13mm, text width=13mm, fill=izpaper] (z) {$\mathbf{z}\in\mathbb{R}^{512}$\\[1pt]{\scriptsize problem\\embedding}};
\node[izhidden, right=of z, minimum height=13mm, text width=22mm] (k) {Kernel embedding\\[1pt]{\scriptsize $\pi_\theta(\kappa\mid\mathbf{z})$ over typed,\\total, emittable kernels}};
\node[izbox, right=of k, minimum height=13mm, text width=17mm, fill=izpaper] (kappa) {Kernel $\kappa$\\[1pt]{\scriptsize frozen before\\any execution}};
\node[izstage, below=8mm of kappa, text width=17mm] (fit) {Fit\\[1pt]{\scriptsize $M=\kappa.\mathrm{fit}(\text{rows})$}};
\node[izstage, left=of fit, text width=22mm] (emit) {Emit\\[1pt]{\scriptsize Rust \textperiodcentered{} Python\\C++17 \textperiodcentered{} JavaScript\textsuperscript{\dag}}};
\node[izstage, left=of emit, text width=22mm] (ver) {Verify\\[1pt]{\scriptsize native and emitted predictions agree (Eq.~\ref{eq:parity})}};
\node[izbox, left=of ver, text width=25mm, draw=izaccent, line width=0.9pt] (src) {Standalone source\\[1pt]{\scriptsize standard library only;\\the dataset's own algorithm}};
\draw[izarrow] (data) -- (enc); \draw[izarrow] (enc) -- (z); \draw[izarrow] (z) -- (k); \draw[izarrow] (k) -- (kappa);
\draw[izarrow] (kappa) -- (fit); \draw[izarrow] (fit) -- (emit); \draw[izarrow] (emit) -- (ver); \draw[izarrow] (ver) -- (src);
\begin{scope}[on background layer]
\node[fit=(k), inner sep=4pt, draw=izink, dashed, rounded corners=3pt, line width=0.5pt] (box) {};
\end{scope}
\node[iznote, above=1.5mm of box] {protected: described by guarantees, not mechanism};
\node[iznote, above=1.5mm of kappa] {IZ-0: one decode};
\end{tikzpicture}

\caption{The data-to-algorithm flow. Only training rows enter; the kernel is frozen before anything
is executed (the track-C contract). The dark stage is the protected kernel embedding. $^\dag$C\texttt{++}17 and
JavaScript emitters cover fewer kernel heads than Rust and Python (\cref{sec:code}).}
\label{fig:pipeline}
\end{figure}

\subsection{The problem encoding}

The encoder reads the values and labels of the training rows and nothing else. Its principal
component is Hyperion v7, a transformer with model width 512, eight attention heads and twelve
blocks over six bands of scale, whose pooled output is a 512-dimensional embedding
(\code{crates/hyperion/src/hyperion\_v7}). The engine's current policy concatenates that embedding
with a 1,024-dimensional foundation encoding and a 512-dimensional encoding of the target column, a
2,048-value context frozen under the recipe \code{frozen-512+1024+512-clip32/v1}
(\code{crates/jope-policy}). Encoders are fitted only on training contexts and are never refitted
when new datasets arrive. The encoder may read training targets --- supervised conditioning is the
point --- but outer-test rows and their labels never reach it.

\subsection{The kernel embedding}

The kernel embedding is a conditional generative model over algorithms. It maps $z$ to a
distribution $\pi_\theta(a\mid z)$ over complete kernels, decoded step by step; the current release
is a decoder of width 512 with six layers, four heads and 27,796,494 parameters
(checkpoint \code{policy-step-60000}). Two qualifications keep this description honest. A
512-dimensional solution embedding $h$ is the bottleneck we propose to test, against direct
conditioning and other widths; neither the 512-dimensional problem signature nor the 2,048-value
policy context establishes that a trained solution embedding of that width exists. And one
proposal is one committed trajectory, not one forward pass: decoding a kernel may take many
autoregressive calls, each counted. A kernel is a feature map followed by a fitted head --- a
ridge regression, a Gaussian process, a gradient-boosted ensemble and so on --- and every kernel the
decoder can express carries six guarantees:
\begin{itemize}
\item \textbf{Typed.} Every operation has an input and output signature; ill-typed compositions
cannot be written.
\item \textbf{Total.} Finite inputs give finite outputs; missing values, constant columns, single
rows and empty feature sets have defined behaviour --- the edge cases the challenge's organizer
checks demand of frontier entrants.
\item \textbf{Bounded.} Work grows at a stated rate in rows and features; there are no unbounded
loops.
\item \textbf{Deterministic.} All randomness flows from one seed.
\item \textbf{Emittable.} Every operation lowers to each target language, or the emitter records a
refusal; a refusal is a failure, never a silent substitution of another language.
\item \textbf{Verifiable.} An outcome counts only when two independent implementations, the native
evaluator and the compiled emission, agree on every prediction:
\begin{equation}
\bigl|\hat y^{\mathrm{native}}_i - \hat y^{\mathrm{emitted}}_i\bigr| \le 10^{-9}\max\bigl(1, |\hat y^{\mathrm{native}}_i|\bigr)
\quad\text{for all rows } i.
\label{eq:parity}
\end{equation}
\end{itemize}

\paragraph{Algorithm decomposition} To learn a distribution over algorithms, an algorithm must first
be decomposed into a canonical sequence of typed operations, the way text is decomposed into
tokens. Decomposition makes algorithms comparable (two programs that compute the same thing receive
the same canonical form), countable (a distribution over them is well defined) and learnable (a
decoder can assign credit to each step). Training the kernel embedding therefore requires a large
corpus of algorithms that have been decomposed this way \emph{and} measured on real data. In this
sense, decomposition is the tokenizer of invention: a model can only learn to propose what has been
broken into parts it can compose. It is also the bridge between our two routes, because a general
algorithm invented by a frontier model can be decomposed into new operations that the kernel
embedding may then use (\cref{sec:protocol}, E10).

\paragraph{What is not yet true} The release policy's proposals have so far been consumed only as
seeds for the engine's evaluation-based proposer, and its first accepted-model run timed out before
producing a proposal (\code{docs/asm.md}). The earlier CLOVER policy, an 893,957-parameter decoder
over a 2,051-value encoding, decodes inside a fixed catalog of 6,640 programs and is therefore
catalog-bound in the sense of \cref{prop:closed} \citep{jopedime2026writeup}. It is a
pre-normalized causal transformer of width 128 with two blocks, four heads, a feed-forward width of
512 and programs of at most 640 representation tokens --- tokens of the kernel language, not of any
natural language. Its outcome head ranks freshly executed programs at chance (pairwise accuracy
0.506, 95\% interval $[0.441,0.568]$, in 32 development contexts), and the repository's archive
records that the neural policy did not beat the non-neural proposer on held-out seed comparisons.
A later offline surrogate raised pairwise ranking accuracy by 0.046 ($[0.040,0.052]$) across 1,136
frozen shadow contexts from 255 parents; that is post-iteration shadow validation, and the
repository keeps its flag \code{flywheel\_established=false}. We report these results because they define the
starting point honestly: \emph{no neural proposer in this system has yet satisfied the track-C
contract on held-out data.}

\subsection{The game}

Each training context $c$ is one dataset cut with two inner folds on its training rows. A move is a
complete kernel. Its outcome is the log ratio of its error to that of the mean predictor fitted on
the complementary fold:
\begin{equation}
q(c,a) = -\tfrac12 \sum_{j=1}^{2} \log \frac{\max(r_j,\varepsilon_c)}{\max(b_j,\varepsilon_c)},
\qquad \varepsilon_c = \max\bigl(10^{-12}\max(1,|\bar y|),\; 10^{-6} s_c\bigr),
\label{eq:q}
\end{equation}
where $r_j$ is the program's RMSE on fold $j$, $b_j$ the mean predictor's, and $s_c$ a robust scale
of the training targets (1.4826 times the median absolute deviation, falling back to the
inter-quartile range and then the standard deviation) \citep{jopedime2026writeup}. $q=0$ is the
mean predictor and $q=\log 10$ a tenfold lower geometric-mean error; $q$ is unit-free, so games on
different datasets are comparable, and it is clipped to $\pm\log 1000$.

The opponent is the \emph{incumbent}: the best verified program found so far for the same context,
its high-water mark $h_c$. A move wins if $q > h_c+\tau$, draws within $\tau=0.01$, and loses
otherwise; any failure, timeout or parity violation loses. Because every incumbent is itself a
program the system once found, we call this \emph{algorithmic self-play}, with one correction to
the analogy: it is not a two-player zero-sum game but a one-step contextual game against a rising
bar, closer to a bandit than to Go. \Cref{fig:game} marks which parts of the loop have run.

\begin{figure}[t]
\centering
% Algorithmic self-play. Solid: implemented and run (first policy generation). Dashed:
% implemented, offline and shadow-only. Dotted: proposed in this paper (gate M7).
\begin{tikzpicture}[x=1cm,y=1cm]
\tikzset{g/.style={izbox, text width=26mm, minimum height=13mm}}
\node[g] (ctx) at (0,0) {Context $c$\\[1pt]{\scriptsize one dataset cut;\\two inner folds}};
\node[izhidden, text width=26mm, minimum height=13mm] (pol) at (3.7,0) {Policy $\pi_\theta$\\[1pt]{\scriptsize proposes kernel $a$}};
\node[izstage, text width=26mm, minimum height=13mm] (exe) at (7.4,0) {Execute $a$\\[1pt]{\scriptsize native and emitted;\\parity-checked}};
\node[g, fill=izpaper] (q) at (11.1,0) {Outcome $q(c,a)$\\[1pt]{\scriptsize \cref{eq:q}}};
\node[g] (teach) at (0,-2.9) {Teachers\\[1pt]{\scriptsize high-water-mark programs\\from an offline search}};
\node[g] (train) at (3.7,-2.9) {Training signal\\[1pt]{\scriptsize imitation, value,\\ranking, feasibility}};
\node[g] (judge) at (7.4,-2.9) {Win, draw or loss\\[1pt]{\scriptsize $q>h_c+\tau$ wins;\\failure loses}};
\node[g] (inc) at (11.1,-2.9) {Incumbent $h_c$\\[1pt]{\scriptsize best verified program\\for $c$ so far}};
\draw[izarrow] (ctx) -- (pol);
\draw[izarrow] (pol) -- (exe);
\draw[izarrow] (exe) -- (q);
\draw[izarrow] (q) -- (inc);
\draw[izarrow] (inc) -- (judge);
\draw[izarrow] (judge) -- (train);
\draw[izarrow] (teach) -- node[iznote, above]{imitation} (train);
\draw[izarrow] (train) -- node[iznote, right, align=left]{update $\theta$} (pol);
\draw[izarrow, dashed] (judge.south) -- ++(0,-0.55) -| node[iznote, pos=0.25, below]{a winner becomes the incumbent (offline shadow pools only)} (inc.south);
\draw[izarrow, dotted, line width=0.9pt] (inc.south east) ++(-0.3,0) -- ++(0,-1.35) -| node[iznote, pos=0.25, below]{verified winners added to the teachers (proposed; gate M7)} (teach.south);
\end{tikzpicture}

\caption{Algorithmic self-play. Solid arrows have run in the first policy generation; the dashed
ratchet runs offline on shadow pools and trains no policy; the dotted path is proposed. Teachers
are verified high-water-mark programs from an offline search, frozen for a model generation;
programs the model writes train its value and ranking heads but never become imitation targets.}
\label{fig:game}
\end{figure}

\begin{table}[t]
\caption{AlphaGo Zero and InnovationZero, row by row. The last two rows are the ones the protocol
in \cref{sec:protocol} is designed to measure.}
\label{tab:go}
\centering\tablesf
\begin{tabularx}{\linewidth}{@{}lLL@{}}
\toprule
& AlphaGo Zero & InnovationZero \\
\midrule
Board & a $19\times19$ position & the training rows of one dataset \\
Board as seen by the network & 17 binary feature planes & a 512-dimensional problem embedding (plus context) \\
Rules & the rules of Go & typed, total, bounded, emittable kernels \\
Move & place one stone & commit one complete algorithm \\
Outcome & win or loss at game end & held-out error relative to the mean predictor, $q$ \\
Opponent & the best previous network & the best previous program for that dataset \\
Policy target & search-improved move distribution & verified incumbents from an offline search \\
Value target & game outcome & measured outcome $q$ \\
Without search & raw network, about 3,000 Elo & track C: pending \\
With search & 5,185 Elo & track S: the engine, $k=256$ \\
What humans supplied & the rules & the operation families and the fundamentals \\
\bottomrule
\end{tabularx}
\end{table}

\subsection{Objectives}

\paragraph{The existing training loss} The first policy generation minimizes a support-masked sum of
six terms \citep{jopedime2026writeup}:
\begin{equation}
\mathcal{L}_6 = 1.00\,\mathcal{L}_{\mathrm{CE}} + 0.25\,\mathcal{L}_{\mathrm{BT}} + 0.10\,\mathcal{L}_{\mathrm{list}}
+ 0.25\,\mathcal{L}_{q} + 0.10\,\mathcal{L}_{\mathrm{feas}} + 0.05\,\mathcal{L}_{\mathrm{rt}} ,
\label{eq:l6}
\end{equation}
where $\mathcal{L}_{\mathrm{CE}}$ is sequence cross-entropy toward the frozen teachers (learn to
write what won), $\mathcal{L}_{q}$ a Huber loss between a predicted and a measured outcome (predict
how well a program will do before running it), $\mathcal{L}_{\mathrm{BT}}$ and
$\mathcal{L}_{\mathrm{list}}$ pairwise and listwise ranking losses on that prediction,
$\mathcal{L}_{\mathrm{feas}}$ a three-way classification of completion, failure and timeout, and
$\mathcal{L}_{\mathrm{rt}}$ a censored loss on log runtime. Continued training adds an exact KL
anchor to the frozen first generation, $\mathcal{L}_{\mathrm{v2}} = \mathcal{L}_6 +
0.2\,\mathrm{KL}(\pi_\theta \,\|\, \pi_{\mathrm{v1}})$. The analogy with AlphaGo Zero is close:
imitation toward search-improved targets supplies the policy, outcome prediction supplies the
value.

\paragraph{The objective track C needs} $\mathcal{L}_6$ has no term that rewards a single proposal
for being good; it imitates teachers and models outcomes. The objective that matches the track-C
contract is the expected utility of one proposal,
\begin{equation}
J_1(\theta) = \mathbb{E}_{c}\,\mathbb{E}_{a\sim\pi_\theta(\cdot\mid z_c)}\bigl[u(c,a)\bigr],
\qquad u = \begin{cases}\mathrm{clip}(q,-1,1) & \text{valid program}\\ -1 & \text{any failure,}\end{cases}
\label{eq:j1}
\end{equation}
optimized, for instance, by a policy gradient with a per-context baseline $b_c$ and a reference
penalty,
\begin{equation}
\mathcal{L}_{\mathrm{PG}} = -\,\mathbb{E}\bigl[(u - \operatorname{sg}(b_c))\log\pi_\theta(a\mid z_c)\bigr]
+ \beta\,\mathbb{E}_c\,\mathrm{KL}\bigl(\pi_\theta \,\|\, \pi_{\mathrm{ref}}\bigr).
\label{eq:pg}
\end{equation}
Two cautions follow. The historical selection statistic, the best-of-eight surplus
\begin{equation*}
J_8=\mathbb{E}\bigl[\max(h_c,q_1,\ldots,q_8)-h_c\bigr],
\end{equation*}
can improve while the first proposal gets worse, and a falling $\mathcal{L}_6$ can too, because its ranking terms train the outcome head rather than the
probability of a complete program; neither is evidence about $J_1$. And if the reference policy gives
zero probability to programs outside its
catalog, $\mathrm{KL}(\pi_\theta\,\|\,\pi_{\mathrm{ref}})$ is infinite for any $\pi_\theta$ that
proposes one, so the existing anchor cannot be carried into an open-support experiment unchanged;
the reference must itself have full support over the legal language. Checkpoints are selected on
$J_1$ at $K=1$, never on best-of-$B$, which can reward a lucky tail without improving any single
proposal.

\paragraph{The objective, in one table} \Cref{tab:objectives} sets the three objectives side by
side. It is the whole ``wrong objective'' argument of \cref{sec:discussion} in miniature.

\begin{table}[h]
\caption{What each kind of system is trained to maximize.}
\label{tab:objectives}
\centering\tablesf
\begin{tabularx}{\linewidth}{@{}lLl@{}}
\toprule
System & Trained to maximize & Ground truth? \\
\midrule
Pretrained language model & likelihood of the next token of human-written text & what a person \emph{would write} \\
Preference-tuned language model & a learned model of human preference & what a person \emph{would prefer} \\
Reasoning-tuned language model & reward from a checker on problems with verifiable answers & what a checker \emph{accepts} \\
InnovationZero (track C) & held-out utility of the one program it commits, $J_1$ & what the code \emph{does} \\
\bottomrule
\end{tabularx}
\end{table}

\subsection{Data, splits and optimization}

The training corpus holds 14,987 data runs from 3,003 dataset versions, of which 7,120 contexts have
complete teacher and input records (6,857 training and 263 validation), drawn from a real tier of
3,859 datasets: 2,611 regression, 843 binary, 322 multi-label and 83 time series
\citep{jopedime2026writeup}. Separation is enforced at five levels. Outer test rows never reach the
encoder, the normalizer, the labels or the teachers. Outcomes are measured on two inner folds of the
training rows. Datasets are grouped into 1,068 parent lineages, about 5\% of which are assigned to
validation by a fixed hash; aliases, target variants and synthetic descendants inherit their
parent's role, so no family straddles training and validation. Sealed families stay outside the
training closure entirely. And the development roster is split into a freely inspected half and a
confirmation half that receives exactly one official look per pre-specified comparison.

The first generation was trained with AdamW (learning rate $3\times10^{-4}$, weight decay 0.01,
gradient clipping at 1, plateau halving with patience 5 to a floor of $10^{-6}$; biases and
normalization weights exempt from decay), drawing up to 16 contexts per update with at most one per
lineage. The continued (v2) run restarts from the first generation's optimizer state and adds the KL
anchor on selected training contexts, so any gain it reports needs restart-only and equal-compute
controls before it can be credited to the anchor. It ran 17,773 updates and plateaued: full
validation loss stayed within $[0.8968, 0.8978]$ for its last 31 epochs. We print the plateau
rather than hide it; it is one more reason the decisive experiment in \cref{sec:protocol} trains on
$J_1$ rather than continuing $\mathcal{L}_6$.
